%% LyX 2.4.4 created this file.  For more info, see https://www.lyx.org/.
%% Do not edit unless you really know what you are doing.
\documentclass[english]{article}
\usepackage[T1]{fontenc}
\usepackage[latin9]{inputenc}
\usepackage{enumitem}
\usepackage{amsmath}
\usepackage{amsthm}
\usepackage{amssymb}
\usepackage{geometry}
\geometry{verbose,tmargin=1in,bmargin=1in,lmargin=1in,rmargin=1in}
\PassOptionsToPackage{normalem}{ulem}
\usepackage{ulem}

\makeatletter
%%%%%%%%%%%%%%%%%%%%%%%%%%%%%% Textclass specific LaTeX commands.
\numberwithin{equation}{section}
\numberwithin{figure}{section}
\newlength{\lyxlabelwidth}      % auxiliary length 

%%%%%%%%%%%%%%%%%%%%%%%%%%%%%% User specified LaTeX commands.
\usepackage{fancyhdr}
\usepackage{lastpage}
\pagestyle{fancy}
\usepackage{enumitem}
\usepackage{ifthen}
\everymath=\expandafter{\the\everymath\displaystyle}
%\DeclareMathSizes{10}{10}{10}{10}
\usepackage{multicol}

\usepackage{tikz}
\usepackage{tikz-3dplot}
\usetikzlibrary{calc,arrows}

\usetikzlibrary{patterns}
\usetikzlibrary{plotmarks}
\usepackage{pgfplots}
\pgfplotsset{compat=newest}

\usepackage{calc}
\usepackage[nomessages]{fp}

\usepackage{datetime}
\usepackage[super]{nth}
% Added by lyx2lyx
\setlength{\parskip}{\smallskipamount}
\setlength{\parindent}{0pt}

\makeatother

\usepackage{babel}
\begin{document}
\renewcommand{\author}{Dr.\,James Ashe}

\newcommand{\classNum}{336}

\newcommand{\classSec}{A}

\newcommand{\examType}{Final Exam Review}

\renewcommand{\date}{\formatdate{25}{11}{2013}}

%%First page's header%%%%%%%%%%%%%%%%%%%%%%
\fancypagestyle{firststyle} %%header settings for first pagr
{
	\fancyhead[L]{MTH \classNum{}(\classSec{}) \author{}\\
	\textbf{\examType{}} ~\date{}}
	\fancyhead[C]{}
	\fancyhead[R]{\textbf{}}
	\setlength{\headsep}{14pt}
}
%%%%%%%%%%%%%%%%%%%%%%%%%%%%%%%%%%%%%%%%%%

%%Next page's header%%%%%%%%%%%%%%%%%%%%%%%
\setlist{nolistsep}
\lhead{\classNum{}(\classSec{})} 
\chead{\textbf{}} 
\cfoot{\thepage{}~of~\pageref{LastPage}} 
\setlength{\headsep}{5pt} 
%%%%%%%%%%%%%%%%%%%%%%%%%%%%%%%%%%%%%%%%%%%

%%Various settings; see the preamble for other settings%%%%%
%\setenumerate[0]{leftmargin=12pt,itemindent=0pt,labelwidth=0pt}
\setlist{topsep=.5pt}
\renewcommand{\labelenumi}{\textbf{\arabic{enumi}.}}
\renewcommand{\labelenumii}{\textbf{\alph{enumii}.}}
\thispagestyle{firststyle}
%%%%%%%%%%%%%%%%%%%%%%%%%%%%%%%%%%%%%%%%%%

\newcommand{\xmax}{0}
\newcommand{\xmin}{-\xmax}
\newcommand{\xstep}{0}
\newcommand{\ymax}{0}
\newcommand{\ymin}{-\ymax}
\newcommand{\ystep}{0} 
\newcommand{\dspace}{\vspace{1cm}}
\newcommand{\xa}{0}
\newcommand{\xb}{0}
\newcommand{\ya}{1}
\newcommand{\yb}{1}

\newcommand{\vectors}[2]{\textbf{#1}_{1},\textbf{#1}_{2},\dots,\textbf{#1}_{#2}}
\newcommand{\vectorsN}[1]{\textbf{#1}_{1},\textbf{#1}_{2},\dots,\textbf{#1}_{n}}
\newcommand{\R}[1]{\mathbb{R}^{#1}}

\textbf{The final date/time: \formatdate{4}{12}{2013} from 10am--12pm.}
\begin{itemize}
\item Let $S={\vectors{v}{m}}$ be a set of $m$ vectors in $\R{n}$.

\begin{itemize}
\item $\mbox{Span}S$ is a subspace of $\R{n}$.

\begin{itemize}
\item $\dim(\mbox{Span}S)\leq m$. $\dim(\mbox{Span}S)$ (the \emph{dimension
}of the span of $S$) is equal to the number of \uline{independent}
vectors in $S$. 
\item Note: $\mbox{Span}S\neq S$, $S$ is not necessarily a subspace (and
usually is not). 
\end{itemize}
\item If $m<n$ then $\mbox{Span}S\neq\R{\mbox{n}}$ (too few vectors).
\item If $m>n$ then $S$ is linearly dependent (more vectors than necessary).
\item If $S$ is linearly independent \uline{and} $m=n$ then $S$ is a
\emph{basis }for $\R{n}$. 
\end{itemize}
\item Let $R$ be the \emph{reduced row echelon form }of a $n\times n$
matrix $A$.

\begin{itemize}
\item If $R$ (or $A$) has a row of zeros then $A$ is singular.
\item If $R$ has a row of zeros then the system of linear equations $A\mathbf{x}=\mathbf{b}$
has infinitely many solutions \uline{or} no solution ($A$ is singular).

\begin{itemize}
\item [\textbf{ex.}]$\begin{bmatrix}\begin{array}{rrr|r}
1 & 0 & 4 & 5\\
0 & 1 & 2 & 3\\
0 & 0 & 0 & 0
\end{array}\end{bmatrix}$ and $\begin{bmatrix}\begin{array}{rrr|r}
1 & 0 & 4 & 0\\
0 & 1 & 2 & 0\\
0 & 0 & 0 & 0
\end{array}\end{bmatrix}$ have infinitely many solutions.
\item [\textbf{ex.}]$\begin{bmatrix}\begin{array}{rrr|r}
1 & 0 & 4 & 5\\
0 & 1 & 2 & 3\\
0 & 0 & 0 & 2
\end{array}\end{bmatrix}$ has no solution.
\end{itemize}
\item If $R$ reduces to the identity then $A\mathbf{x}=\mathbf{b}$ has
a \uline{unique} solution ($A$ is nonsingular).

\begin{itemize}
\item [\textbf{ex.}]$\begin{bmatrix}\begin{array}{rrr|r}
1 & 0 & 0 & 5\\
0 & 1 & 0 & 0\\
0 & 0 & 1 & 1
\end{array}\end{bmatrix}$ and $\begin{bmatrix}\begin{array}{rrr|r}
1 & 0 & 0 & 0\\
0 & 1 & 0 & 0\\
0 & 0 & 1 & 0
\end{array}\end{bmatrix}$ have unique solutions (a unique solution to $A\mathbf{x}=\mathbf{0}$
is called a \emph{trivial} solution).
\end{itemize}
\end{itemize}
\end{itemize}
\begin{multicols}{2}
\begin{itemize}
\item Let $A$ be an $n\times n$ matrix then $A$ is \emph{nonsingular}

\begin{itemize}
\item [$\iff$] $A$ reduces to the identity matrix
\item [$\iff$] $A^{-1}$ exists ($A$ is invertible)
\item [$\iff$]$A\mathbf{x}=\mathbf{b}$ has a unique solution.
\item [$\iff$]$A\mathbf{x}=\mathbf{0}$ has a \emph{trivial solution} ($\mathbf{x}=\mathbf{0}$).
\item [$\iff$] $\det A\neq0$\columnbreak
\end{itemize}
\item Let $A$ be an $n\times n$ matrix then $A$ is \emph{singular}

\begin{itemize}
\item [$\iff$] $A$ reduces to a matrix with a row of zeros.
\item [$\iff$] $A^{-1}$ does not exists.
\item [$\iff$]$A\mathbf{x}=\mathbf{b}$ has no or infinite solutions.
\item [$\iff$]$A\mathbf{x}=\mathbf{0}$ has a \emph{non-trivial solution}
($\mathbf{x}\neq\mathbf{0}$).
\item [$\iff$] $\det A=0$
\end{itemize}
\end{itemize}
\end{multicols}
\begin{itemize}
\item Let $A$ and $B$ be $n\times n$ matrices then

\begin{itemize}
\item $\det(AB)=\det(A)\det(B)$
\item $\frac{1}{\det A}=\det(A^{-1})$ $\iff$ $\det A\neq0$ ($A^{-1}$
is invertible)
\item for an $A'$ obtained from $A$, $\det A=\lambda\det(A')$ for some
nonzero $\lambda\in\mathbb{R}$ where the value of $\lambda$ can
be found by using the properties of determinants.
\item $\det A=\det B$ if $A$ and $B$ are similar. \newpage
\end{itemize}
\item Eigenvalues and Eigenvectors.

\begin{itemize}
\item $p_{A}(t)=\det(A-\lambda I_{n})=0$ $\iff$ $\lambda$ is an eigenvalue.
\item $A$ and $B$ are \emph{similar} \emph{matrices} if $A=P^{-1}BP$.

\begin{itemize}
\item $\Rightarrow$ $\det A=\det B$. 
\item $\Rightarrow$ $p_{A}(t)=p_{A}(t)$, i.e., similar matrices have the
same eigenvalues.
\end{itemize}
\item $0$ is an eigenvalue of $A$ $\iff$$A$ is singular.
\item If $A$ has distinct eigenvalues $\lambda_{1},\lambda_{2},\dots,\lambda_{n}$
with corresponding eigenvectors $S=\vectors{v}{n}$, then $S$ is
linearly independent and $A$ is diagonalizable (this is a sufficient
but not necessary condition).
\end{itemize}
\item Linear transformations, $T\colon\R{\mbox{m}}\rightarrow\R{\mbox{n}}$.

\begin{itemize}
\item If $T$ is a linear transformation then there exists a \emph{standard
}$m\times n$ \emph{matrix }$A$ such that $\mu_{A}\colon\R{\mbox{m}}\rightarrow\R{\mbox{n}}=T$.
\item $T[\mathbf{x}]=\mathbf{y}$ is a function given by $A\mathbf{x}=\mathbf{y}$
where $\mathbf{x}\in\R{\mbox{m}}$ and $\mathbf{y}\in\R{\mbox{n}}$.
\item The $i^{\mbox{th}}$ column of a standard matrix will be $T\left[\mathbf{x}_{i}\right]$
where $\mathbf{x}_{i}$ is the $i^{\mbox{th}}$vector in the standard
basis for $\R{\mbox{m}}$.
\end{itemize}
\item Be able to compute the following:

\begin{itemize}
\item $\det A$ using cofactors and properties of matrices.
\item $A^{-1}$.

\begin{itemize}
\item If $A=\begin{bmatrix}\begin{array}{rr}
a & b\\
c & d
\end{array}\end{bmatrix}$ and $\det A\neq0$ then $A^{-1}=\begin{bmatrix}\begin{array}{rr}
\frac{d}{ad-bc} & \frac{-b}{ad-bc}\\
\frac{-c}{ad-bc} & \frac{a}{ad-bc}
\end{array}\end{bmatrix}$.
\item If $\det A\neq0$ and $\det B\neq0$ then $\left(\begin{bmatrix}\begin{array}{rr}
A & \mathbf{0}\\
\mathbf{0} & B
\end{array}\end{bmatrix}\right)^{-1}=\begin{bmatrix}A^{-1} & \mathbf{0}\\
\mathbf{0} & B^{-1}
\end{bmatrix}$ (see ex 12 in section 2.3)
\end{itemize}
\item $p_{A}(t)$.
\item the reduced row echelon form of $A$.
\item the eigenvalues and eigenvectors for a given matrix.
\end{itemize}
\end{itemize}

\end{document}
